\documentclass[tikz,border=8pt]{standalone}
\usepackage{tikz}
\usepackage{amsmath}
\usepackage{pgfplots}
\pgfplotsset{compat=1.18}
\usepackage{ctex}
\begin{document}
\begin{tikzpicture}
\begin{axis}[
  width=10cm, height=8cm,
  xlabel={$x_1$},
  ylabel={$x_2$},
  zlabel={$f(x_1,x_2)$},
  title={二元标准正态 PDF（$\rho=0$，独立）},
  title style={font=\small\bfseries},
  xlabel style={font=\small},
  ylabel style={font=\small},
  zlabel style={font=\small},
  xmin=-3, xmax=3,
  ymin=-3, ymax=3,
  zmin=0, zmax=0.17,
  view={40}{30},
  colormap/cool,
  xticklabel style={font=\tiny},
  yticklabel style={font=\tiny},
  zticklabel style={font=\tiny},
  samples=30,
  clip=false,
]

\addplot3[surf, shader=faceted, opacity=0.85, domain=-3:3]
  {exp(-(x^2+y^2)/2)/(2*pi)};

\end{axis}
\node[font=\scriptsize, fill=yellow!15, draw, rounded corners=2pt, text width=6cm]
  at (5,-0.5)
  {$\rho=0$：等值面为圆形，$X_1\perp X_2$\\
   PDF 最高点在均值 $(\mu_1,\mu_2)=(0,0)$ 处};
\end{tikzpicture}
\end{document}
